\section{Method} \label{method}
In this section, we first introduce the Dirichlet Multinomial Mixture (DMM) model~\cite{nigam2000text}. Next, we provide a detailed derivation and analysis of the proposed GSDMM. Finally, we identify the limitations of GSDMM and present an enhanced model-based approach for short text clustering to address these issues.

\subsection{The DMM Model}

\begin{table}
\caption{Notations}
%\begin{small}
\centering
%\setlength{\abovecaptionskip}{-0.002cm}
%\setlength{\belowcaptionskip}{-0.28cm}
%\small\addtolength{\tabcolsep}{-4pt}
\begin{tabular}{cl}
  \hline
  % after \\: \hline or \cline{col1-col2} \cline{col3-col4} ...
  $V$ & size of the vocabulary \\
  $D$ & number of documents in the corpus \\
  $I$ & number of iterations \\
  $\bar{L}$ & average length of documents \\
  $\vec{d}$  & documents in the corpus\\
  $\vec{z}$  & cluster assignments of each document \\
  $m_z$ & number of documents in cluster $z$ \\
  $n_z$ & number of words in cluster $z$ \\
  $n_z^{w}$   & frequency of word $w$ in cluster $z$ \\
  $N_d$   & number of words in document $d$ \\
  $N_d^{w}$   & frequency of word $w$ in document $d$ \\
  $\alpha$ & pseudo number of documents in each cluster \\
  $\beta$ & pseudo frequency of each word in each cluster \\
%  &  \\
  \hline
\end{tabular}
\label{tab:Notations}
%\end{small}
\end{table}

The DMM model~\cite{nigam2000text} is a probabilistic generative model for documents that is based on two key assumptions: 1) the documents are generated by a mixture model~\cite{mclachlan_basford88}, and 2) there is a one-to-one correspondence between the mixture components and the clusters. The generative process of the DMM model is illustrated as follows:
\begin{align}
    & \label{gen:dmm1}
    \theta|\alpha \sim Dir(\alpha) \\
    & \label{gen:dmm2}
    z_d|\theta \sim Mult(\theta) \qquad d=1,...,D\\
    & \label{gen:dmm3}
    \phi_k|\beta \sim Dir(\beta) \qquad k=1,...,K\\
    & \label{gen:dmm4}
    d|z_d,\{\phi_k\}_{k=1}^K \sim p(d | \phi_{z_d})
\end{align}
Here, ``$X \sim S$'' means ``$X$ is distributed according to $S$'', so the right side is a specification of distribution. When generating document $d$, the DMM model first selects a mixture component (cluster) $z_d$ for document $d$ according to the mixture weights (weights of clusters), $\theta$, in Equation~\eqref{gen:dmm2}. Then, document $d$ is generated by the selected mixture component (cluster) from distribution $p(d | \phi_{z_d})$ in Equation~\eqref{gen:dmm4}. The weight vector of the clusters, $\theta$, is generated by a Dirichlet distribution with a hyper-parameter $\alpha$, as in Equation~\eqref{gen:dmm1}. The cluster parameters $\phi_z$ are also generated by a Dirichlet distribution with a hyper-parameter $\beta$, as in Equation~\eqref{gen:dmm3}.

The probability of document $d$ generated by cluster $z_d$ is defined as follows:
\begin{equation} \label{eq:Pdz}
    p(d | \phi_{z_d}) = \prod_{w \in d} Mult(w|\phi_{z_d})
\end{equation}
Here, we make the Naive Bayes assumption: the words in a document are generated independently when the document's cluster assignment $z_d$ is known. We also assume that the probability of a word is independent of its position within the document.

\subsection{The GSDMM Algorithm} \label{GSDMM}
In this part, we give the derivation of the collapsed Gibbs sampling algorithm for the DMM model. The documents $\vec{d}=\{d_i\}_{i=1}^D$ are observed and the cluster assignments $\vec{z} = \{z_i\}_{i=1}^D$ are latent. The detail of our GSDMM algorithm is shown in Algorithm~\ref{Algo:GSDMM}, and the meaning of its variables is shown in Table~\ref{tab:Notations}. The conditional probability of document $d$ choosing cluster $z$ given the information of other documents and their cluster assignments can be factorized as follows:
\begin{align}
    & \notag
    p(z_d=z | \vec{z}_{\neg{d}}, \vec{d}, \alpha, \beta)\\
    & \label{eq:dpm1.1}
    \propto p(z_d=z | \vec{z}_{\neg{d}}, \vec{d}_{\neg{d}}, \alpha, \beta)
            p(d | z_d=z, \vec{z}_{\neg{d}}, \vec{d}_{\neg{d}}, \alpha, \beta)\\
    & \label{eq:dpm1.2}
    \propto p(z_d=z | \vec{z}_{\neg{d}}, \alpha)
            p(d | z_d=z, \vec{d}_{z, \neg{d}}, \beta)
\end{align}
Here, we use the Bayes Rule in Equation~\eqref{eq:dpm1.1} and apply the properties of D-Separation~\cite{bishop2006pattern} in Equation~\eqref{eq:dpm1.2}.

The first term in Equation~\eqref{eq:dpm1.2} indicates the probability of document $d$ choosing cluster $z$ when we know the cluster assignments of other documents. The second term in Equation~\eqref{eq:dpm1.2} can be considered as a predictive probability of document $d$ given $\vec{d}_{z, \neg{d}}$, i.e., the other documents currently assigned to cluster $z$.

The first term in Equation~\eqref{eq:dpm1.2} can be derived as follows:
%\begin{small}%equation
\begin{align}
    & \notag
    p(z_d=z | \vec{z}_{\neg{d}}, \alpha)\\
    & \label{eq:dpm2.1}
    = \int{p(z_d=z, \theta | \vec{z}_{\neg{d}}, \alpha)
           d\theta}\\
    & \label{eq:dpm2.2}
    = \int{p(\theta | \vec{z}_{\neg{d}}, \alpha)
           p(z_d=z | \vec{z}_{\neg{d}}, \theta, \alpha)
           d\theta}\\
    & \label{eq:dpm2.3}
    = \int{p(\theta | \vec{z}_{\neg{d}}, \alpha)
           p(z_d=z | \theta)
           d\theta}\\
    &\label{eq:dpm2.4.1}
    = \int{Dir(\theta | \vec{m}_{\neg{d}} + \vec{\alpha})
           Mult(z_d=z | \theta)
           d\theta}\\
    &\label{eq:dpm2.4.2}
    = \int{ \frac{1}{\Delta(\vec{m}_{\neg{d}}+\vec{\alpha})}
        \theta_z \prod_{k=1,k\neq z}^K \theta_k^{m_{k,\neg{d}}+\alpha-1}
         d\theta}\\
    &\label{eq:dpm2.4.3}
    = \frac{\Delta(\vec{m} + \vec{\alpha})}
           {\Delta(\vec{m}_{\neg{d}} + \vec{\alpha})}\\
    &\label{eq:dpm2.4.4}
    = \frac{\prod_{k=1}^K \Gamma(m_k + \alpha)}
            {\Gamma(\sum_{k=1}^K (m_k + \alpha))}
      \frac{\Gamma(\sum_{k=1}^K (m_{k,\neg{d}} + \alpha))}
            {\prod_{k=1}^K \Gamma(m_{k,\neg{d}} + \alpha)} \\
    &\label{eq:dpm2.4.5}
    = \frac{\Gamma(m_{z,\neg{d}} + \alpha + 1)}
            {\Gamma(m_{z,\neg{d}} + \alpha)}
      \frac{\Gamma(D - 1 + K\alpha)}
            {\Gamma(D + K\alpha)} \\
    &\label{eq:dpm2.4.6}
    = \frac{m_{z,{\neg{d}}} + \alpha}
           {D - 1 + K\alpha}
\end{align}
%\end{small}%equation
Here, Equation~\eqref{eq:dpm2.1} exploits the sum rule of probability~\cite{bishop2006pattern}. We use the product rule of probability~\cite{bishop2006pattern} in Equation~\eqref{eq:dpm2.2} and apply the properties of D-Separation in Equation~\eqref{eq:dpm2.3}. The posterior distribution of $\theta$ is a Dirichlet distribution because we assume Dirichlet prior $Dir(\theta | \alpha)$ for the Multinomial distribution $Mult(\vec{z} | \theta)$. In Equation~\eqref{eq:dpm2.4.3}, we adopt the $\Delta$ function used in~\cite{Hei09}, which is defined as $\Delta(\vec{\alpha}) = \frac{\prod_{k=1}^K \Gamma(\alpha)}{\Gamma(\sum_{k=1}^K \alpha)}$. Using the property of $\Gamma$ function: $\Gamma(x+1)=x\Gamma(x)$, we can get Equation~\eqref{eq:dpm2.4.6} from Equation~\eqref{eq:dpm2.4.5}. In Equation~\eqref{eq:dpm2.4.6}, $m_{z,{\neg{d}}}$ is the number of documents in cluster $z$ without considering document $d$, and $D$ is the total number of documents in the dataset. Equation~\eqref{eq:dpm2.4.6} indicates that document $d$ will tend to choose larger clusters when we only consider the cluster assignments of the other documents.

\begin{algorithm}
\SetAlFnt{\small}
%\setlength{\abovecaptionskip}{-0.02cm}
%\setlength{\belowcaptionskip}{-1.5cm}
\DontPrintSemicolon
\KwData{Documents in the input, $\vec{d}$. }
\KwResult{Cluster labels of each document, $\vec{z}$.}
\Begin{
    Initialize $m_z, n_z$, and $n_z^w$ as zero for each cluster $z$\;
    \For{each document $d\in [1,D]$}{
        Sample a cluster for $d$: $z_d \leftarrow z \sim Multinomial(1/K)$\;
        $m_z \leftarrow m_z+1$ and $n_z \leftarrow n_z+N_d$\;
        \For{each word $w\in d$}{
            $n_z^w \leftarrow n_z^w +N_d^w$\;
        }
    }
    \For{$i \in [1,I]$}{
        \For{each document $d\in [1,D]$}{
            Record the current cluster of $d$: $z=z_d$\;
            $m_z \leftarrow m_z-1 $ and $ n_z \leftarrow n_z-N_d$\;
            \For{each word $w \in d$}{
                $n_z^w \leftarrow n_z^w- N_d^w$\;
            }
            Sample a cluster for $d$: $z_d \leftarrow z \sim p(z_d=z|\vec{z}_{\neg{d}}, \vec{d})$ (Equation~\eqref{eq:dpm3.4.7})\;
            $m_z \leftarrow m_z+1 $ and $ n_z \leftarrow n_z+N_d$\;
            \For{each word $w \in d$}{
                $n_z^w \leftarrow n_z^w+ N_d^w$\;
            }
        }
    }
}
\caption{GSDMM}
\label{Algo:GSDMM}
\end{algorithm}

Then, we derive the second term in Equation~\eqref{eq:dpm1.2} as follows:
\begin{align}
    & \notag
    p(d | z_d=z, \vec{d}_{z, \neg{d}}, \beta)\\
    & \label{eq:dpm3.1}
    = \int{p(d, \phi_z | z_d=z, \vec{d}_{z, \neg{d}}, \beta)d\phi_z}\\
    & \label{eq:dpm3.2}
    = \int{p(\phi_z | z_d=z, \vec{d}_{z, \neg{d}}, \beta)
           p(d | \phi_z, z_d=z, \vec{d}_{z, \neg{d}}, \beta)
           d\phi_z}\\
    & \label{eq:dpm3.3}
    = \int{p(\phi_z | \vec{d}_{z, \neg{d}}, \beta)
           p(d | \phi_z, z_d=z)
           d\phi_z}\\
    & \label{eq:dpm3.4.1}
    = \int{Dir(\phi_z | \vec{n}_{z,\neg{d}} + \beta)
           \prod_{w \in d} Mult(w | \phi_z)
           d\phi_z}\\
    & \label{eq:dpm3.4.2}
    = \int{\frac{1}{\Delta(\vec{n}_{z,\neg{d}} + \beta)}
            \prod_{t=1}^V \phi_{z,t}^{n_{z,\neg{d}}^t + \beta - 1}
           \prod_{w \in d} \phi_{z,w}^{n_d^w}
            d\phi_z}\\
    & \label{eq:dpm3.4.3}
    = \frac{\Delta(\vec{n}_z + \beta)}
           {\Delta(\vec{n}_{z,\neg{d}} + \beta)}\\
    & \label{eq:dpm3.4.4}
    = \frac{\prod_{t=1}^V \Gamma(n_z^t + \beta)}
            {\Gamma(\sum_{t=1}^V (n_z^t + \beta))}
      \frac{\Gamma(\sum_{t=1}^V (n_{z,\neg{d}}^t + \beta))}
            {\prod_{t=1}^V \Gamma(n_{z,\neg{d}}^t + \beta)} \\
    & \label{eq:dpm3.4.5}
    =  \frac{\prod_{w \in d} \prod_{j=1}^{N_d^w} (n_{z,\neg{d}}^w + \beta + j - 1)}
           {\prod_{i=1}^{N_d} (n_{z,\neg{d}} + V\beta + i - 1)}
\end{align}
Here, Equation~\eqref{eq:dpm3.1} exploits the sum rule of probability~\cite{bishop2006pattern}. We use the product rule of probability in Equation~\eqref{eq:dpm3.2} and apply the properties of D-Separation \cite{bishop2006pattern} to obtain Equation \eqref{eq:dpm3.3}.
The posterior distribution of $\phi_z$ is a Dirichlet distribution because we assume Dirichlet prior $Dir(\phi_z | \beta)$ for the Multinomial distribution $Mult(\vec{d}_z | \phi_z)$. Because the $\Gamma$ function has the following property: $\frac{\Gamma(x+m)}{\Gamma(x)} = \prod_{i=1}^m(x+i-1)$, we can get Equation \eqref{eq:dpm3.4.5} from Equation~\eqref{eq:dpm3.4.4}.
In Equation~\eqref{eq:dpm3.4.5}, $N_d^w$ and $N_d$ are the number of occurrences of word $w$ in document $d$ and the total number of words in document $d$, respectively, and $N_d = \sum_{w \in d} N_d^w$.
Besides, $n_{z,\neg{d}}^w$ and $n_{z,\neg{d}}$ are the number of occurrences of word $w$ in cluster $z$ and the total number of words in cluster $z$ without considering document $d$, respectively, and $n_{z,\neg{d}} = \sum_{w=1}^V n_{z,\neg{d}}^w$.
We can notice that Equation~\eqref{eq:dpm3.4.5} evaluates some similarity between document $d$ and cluster $z$, and document $d$ will tend to choose a cluster whose documents share more words with it.

Finally, we have the probability of document $d$ choosing cluster $z_d$ given the information of other documents and their cluster assignments as follows:
\begin{align}
& \notag
p(z_d=z | \vec{z}_{\neg{d}}, \vec{d}, \alpha, \beta)\\
&
\label{eq:dpm3.4.7}
\propto \frac{m_{z,\neg{d}} + \alpha}{D-1 + K\alpha}
        \frac{\prod_{w \in d} \prod_{j=1}^{N_d^w} (n_{z,\neg{d}}^w + \beta + j - 1)}
           {\prod_{i=1}^{N_d} (n_{z,\neg{d}} + V\beta + i - 1)}
\end{align}
The first part of Equation~\eqref{eq:dpm3.4.7} relates to Rule 1 of GSDMM (i.e., choose a cluster with more documents). This is also known as the ``richer get richer" property, leading larger clusters to get larger~\cite{teh2010dirichlet}. The second part of Equation~\eqref{eq:dpm3.4.7} relates to Rule 2 of GSDMM (i.e., choose a cluster that is more similar to the current document), which defines the similarity between the document and the cluster. Notably, GSDMM represents each document by its words and their respective frequencies. It models each cluster as a single large document composed of all documents within the cluster while recording the cluster size.

\subsection{The \sysname{} Algorithm} \label{VEMC}
Although the proposed GSDMM demonstrates outstanding performance, an in-depth analysis reveals several aspects that warrant further refinement: 1) Due to the lack of initial document information, random initialization leads to poor performance in early iterations and requires more iterations to correct misallocated documents. 2) When assigning documents to clusters, the uniform word weighting scheme ignores the varying contributions of different words to clusters, potentially introducing noise. 3) GSDMM assumes that at most $K_{max}$ clusters are in the corpus, which is a weaker assumption than specifying the true number of categories. Although GSDMM can prune redundant clusters within a limited number of iterations, it still risks producing a suboptimal number of clusters.

To address these issues, we propose an improved short text clustering method \sysname{} which is shown in Algorithm~\ref{Algo:VEMC}. First, we replace random initialization with adaptive clustering initialization to accelerate model convergence. Next, we incorporate entropy to measure word importance and dynamically adjust word weights, enabling the model to focus more on key terms. Finally, we dynamically refine the number of clusters to better align the predicted distribution with the true category distribution. In this section, we provide a detailed explanation of each improvement.

\begin{algorithm}[t]
%\setlength{\abovecaptionskip}{-0.002cm}
\setlength{\belowcaptionskip}{-0.10cm}
%\begin{small}
\SetAlFnt{\scriptsize}
%%\setlength{\abovecaptionskip}{-0.002cm}
%\setlength{\belowcaptionskip}{-1.5cm}
\DontPrintSemicolon
\KwData{Document vector $\vec{d}$, Maximum number of clusters $K_{max}$, Real number of clusters $K_{real}$.}
\KwResult{cluster assignments of each document $\vec{z}$.}
\Begin{
    //Adaptive Clustering Initialization\;
    Obtain the initial cluster assignments for documents by Algorithm \ref{Algo:ACI}.\;
    //Collapsed Gibbs Sampling with Entropy\;
    \For{$i \in [1,I]$}{
        \For{each document $d\in [1,D]$}{
            Record the current cluster of $d$: $z=z_d$\;
            $m_z \leftarrow m_z-1 $ and $ n_z \leftarrow n_z-N_d$\;
            \For{each word $w \in d$}{
                $n_z^w \leftarrow n_z^w- N_d^w$\;
            }
            \If{$n_z == 0$}{
                //Remove the empty cluster\;
                $K \leftarrow K-1$\;
                Re-arrange cluster indices so that $1,...,K$ are active (i.e., non-empty);\;
            }
            \If{d \% 15 == 0}{
                Update the word entropy values.\;
            }
            Compute the probability of document $d$ and choose a new $z$ from each of the $K$ existing clusters with Equation \eqref{DMM-Entropy}.\;
            $z_d \leftarrow z$\;
            $m_z \leftarrow m_z+1 $ and $ n_z \leftarrow n_z+N_d$\;
            \For{each word $w \in d$}{
                $n_z^w \leftarrow n_z^w+ N_d^w$\;
            }
        }
    }
    //Granularity Adjustment with Cluster Merging\;
    \If{$K > K_{\text{real}}$}{
        Compute the TF-ICF word distributions of clusters by Equation \eqref{idf} and Equation \eqref{tfidf}.\;
        Build a priority queue $Q$ of cluster pairs.\;
        \While{$K > K_{\text{real}}$}{
            //Highest-similarity pair\;
            $(c_p, c_q) \leftarrow \text{PopMax}(Q)$\;
            $m_{c_m} \leftarrow m_{c_p} + m_{c_q}$\;
            $n_{c_m} \leftarrow n_{c_p} + n_{c_q}$\;
            \For{each word $w$}{
                $n_{c_m}^w \leftarrow n_{c_p}^w + n_{c_q}^w$\;
            }
            $K \leftarrow K - 1$\;
            \For{each active cluster $c_a$}{
                Compute the cosine similarity of $(c_m, c_a)$ and insert it into $Q$.\;
            }
        }
    }
}
\caption{\sysname{}}
\label{Algo:VEMC}
%\end{small}
\end{algorithm}

\subsubsection{Adaptive Clustering Initialization}
Traditional random initialization assigns a random label to each document, often ignoring distance information between documents. This results in poor initial performance and requires more iterations to achieve stability. Therefore, we introduce a new initialization method, presented in Algorithm~\ref{Algo:ACI}, which incorporates the similarity of the document cluster in the initialization stage to improve the effectiveness.

% \begin{align}
% & \notag
% p(z_d=z | \vec{z}, \vec{d}, \alpha, \beta)\\
% &
% \label{init}
% \propto \frac{m_{z} + \alpha}{D-1 + K\alpha}
%         \frac{\prod_{w \in d} \prod_{j=1}^{N_d^w} (n_{z}^w + \beta + j - 1)}
%            {\prod_{i=1}^{N_d} (n_{z} + V\beta + i - 1)}
% \end{align}

The document set is initially composed of $K$ randomly selected documents, each serving as an initial cluster. Unlike K-means and the standard Voronoi diagram~\cite{aurenhammer2000voronoi}, which assign each data point to the nearest fixed cluster center, our method iteratively refines clusters during initialization. Each document is assigned to its most similar cluster, and the clusters are continuously updated based on the assigned documents. This process provides a more interpretable and stable starting point for subsequent clustering optimization.

% We draw inspiration from the Voronoi diagram \cite{aurenhammer2000voronoi} to design a novel initialization method, termed Adaptive clustering Initialization (AKI). Unlike K-means and standard Voronoi diagram \cite{aurenhammer2000voronoi}, which assigns each data point to the nearest fixed cluster center, AKI incorporates a dynamic adjustment mechanism for cluster centers during the initialization phase. This adaptive process enhances the alignment between cluster centers and the underlying data distribution, thereby improving the initial partitioning's robustness and semantic coherence.
% Based on the geometric properties of Voronoi diagrams, the $K$ selected documents can be considered as centroids or "seed points." Dividing the space based on the distances to these centroids generates corresponding Voronoi cells. In the context of text classification, these points can be interpreted as the $K$ representative documents, where each document is assigned to the most appropriate cluster center.

\subsubsection{Entropy-based Word Weighting}
The standard Dirichlet multinomial mixture model used in GSDMM assumes an uniform Dirichlet prior parameter ($\beta$) for all words. This assumption implies that all words contribute equally to cluster formation, oversimplifying the nature of textual data. In reality, words exhibit varying levels of discriminative power across categories. 

To address this issue, we introduce entropy-based word weighting to adaptively adjust the influence of individual words based on their informational significance within each cluster.
For each word $w \in V$, its entropy $H(w)$ across the $K$ clusters is computed as follows:
\begin{align}
& \label{pk}
p_k(w) = \frac{n_k^w+\epsilon}{\sum_{k=1}^K n_k^w+K\epsilon} \\
& \label{hw}
H(w) = - \sum_{k=1}^K p_k(w) \log p_k(w)
\end{align}
where $p_k(w)$ represents the proportion of word $w$ in cluster $k$ relative to its total occurrences across all clusters, estimated based on the current cluster assignments. A small value $\epsilon$ is used to prevent division by zero errors. The entropy $H(w)$ measures the uncertainty in the word's distribution across clusters. Words with lower entropy are more informative, as they are predominantly associated with specific clusters. Conversely, words with higher entropy are more evenly distributed across clusters, making them less useful for clustering.

To incorporate word weighting based on entropy, we utilize the entropy of each word to replace the hyper-parameter $\beta$, which adjusts the influence of each word. This approach allows us to compute the probability of document $d$ selecting cluster $z$ as follows:
\begin{align}
& \notag
p(z_d=z | \vec{z}_{\neg{d}}, \vec{d}, \alpha)\\
&
\label{DMM-Entropy}
\propto \frac{m_{z,\neg{d}} + \alpha}{D-1 + K\alpha}
        \frac{\prod_{w \in d} \prod_{j=1}^{N_d^w} (n_{z,\neg{d}}^w + H(w) + j - 1)}
           {\prod_{i=1}^{N_d} (n_{z,\neg{d}} + \sum_{w \in V} H(w) + i - 1)}
\end{align}

This adjustment ensures that Gibbs sampling accounts for the relative importance of each word based on its entropy, allowing the model to prioritize words with higher discriminative power during cluster assignment. Specifically, in Equation~\eqref{DMM-Entropy}, words with higher entropy reduce the influence of the word's frequency in cluster $z$, thereby diminishing their impact on clustering. In contrast, words with lower entropy increase the effect of the frequency $n_{z,\neg{d}}^w$, enhancing their distinguishing power.

This approach enhances the influence of discriminative words while mitigating the noise, allowing for a greater focus on meaningful terms. By prioritizing words that contribute higher mutual information, it aligns with information-theoretic principles and strengthens the model's ability to distinguish clusters. Furthermore, by dynamically adjusting word weights based on entropy, \sysname{} significantly improves upon GSDMM and enhances clustering accuracy, particularly for complex textual datasets.

\begin{algorithm}[t]
\setlength{\belowcaptionskip}{-0.10cm}
\SetAlFnt{\scriptsize}
\DontPrintSemicolon
\KwData{Document vector $\vec{d}$, maximum number of clusters $K_{max}$.}
\KwResult{Initial cluster assignments of each document $\vec{z}$.}
\Begin{
    \For{each cluster $z \in [1, K_{max}]$}{
        Sample a document $d$ with a uniform distribution.\;
        $z_d \leftarrow z$\;
        $m_z \leftarrow m_z+1 $ and $ n_z \leftarrow n_z+N_d$\;
        \For{each word $w \in d$}{
            $n_z^w \leftarrow n_z^w+ N_d^w$\;
        }}
    Construct an empty document set $S$. \;
    Add the sampled documents into $S$. \;
    \For{each document $d \in [1, D]$}{
        \If{$d$ is not in $S$} {
            Sample a new cluster $z$ for $d$: \;  
            $z_d \leftarrow z \sim p(z_d = z| \vec{z}, \vec{d}, \alpha, \beta)$ \;($\vec{d}$ represents the documents and $\vec{z}$ denotes the corresponding cluster assignments in $S$.) \;
            Add the document $d$ into $S$. \;
            $m_z \leftarrow m_z + 1$ and $n_z \leftarrow n_z + N_d$\;
            \For{each word $w \in d$}{
                $n_z^w \leftarrow n_z^w + N_d^w$\;
            }
        }
    }
}
\caption{Adaptive Clustering Initialization}
\label{Algo:ACI}
\end{algorithm}

\subsubsection{Granularity Adjustment with Cluster Merging}
In the previous modules, we generate a fine-grained cluster distribution, also known as high granularity. This process produces more small clusters, capturing subtle local patterns during the iterations. We dynamically adjust clustering granularity through cluster merging to better align the predicted distribution with the true category distribution, effectively consolidating redundant clusters and enhancing clustering performance. In this section, we design a modified Term Frequency-Inverse Document Frequency (TF-IDF) measure, called Term Frequency-Inverse Cluster Frequency (TF-ICF), which uses the weight $W_{t,c}$ to represent a term’s importance to a topic $c$. 
% The clustering performance of GSDMM is highly sensitive to the hyper-parameter $\beta$. Our experiments indicate that this sensitivity arises because the final number of clusters strongly depends on $\beta$, making it difficult to align the number of clusters with the true data categories and challenging to determine an optimal $\beta$. 
\begin{align}
&\label{idf}
    icf_{t} = 1 + \log \left(\frac{1 + K}{1 + cf(t)} \right) \\
&\label{tfidf}
    W_{t,c} = tf_{t,c} \cdot icf_{t}
% &\label{l2norm}
%     W_{c}=\frac{W_{c}^{-1}}{\|W_{c}^{-1}\|_{2}}
% l2norm怎么在正文中体现，需要体现吗
\end{align}
where $K$ is the total number of clusters, term $t$ represents a word, and the cluster frequency $cf(t)$ is the number of clusters that contain the term $t$. The inverse cluster frequency ($icf_{t}$) measures how much information a term $t$ provides to a cluster. Then, the term frequency ($tf_{t,c}$) models the frequency of term $t$ in a cluster $c$, and the cluster $c$ is a single large document composed of all documents within the cluster while recording the cluster size. TF-ICF models the importance of words in clusters instead of individual documents, which allows us to generate topic-word distributions for each cluster of documents. 

Central to our method is the quantification of cosine similarity between clusters based on their topic-word distributions, making it well-suited for comparing semantic similarity. Specifically, to optimize the merging process and reduce computational overhead, we employ a priority queue that stores cluster pairs sorted by similarity scores. Finally, by iteratively merging the TF-ICF representations of the most semantically coherent clusters, we can reduce the number of topics to a user-specified value. This approach brings the resulting number of clusters closer to the true number of categories, yielding optimal results.
% Over-clustering often arises due to the high dimensionality and sparsity inherent in short text data, where minor variations in word distributions can lead to the formation of distinct clusters that do not correspond to meaningful semantic distinctions. SCM mitigates this by leveraging the similarity of word distributions across clusters to identify and merge redundant or semantically similar clusters. This approach not only reduces the number of clusters to a more manageable and interpretable level but also enhances the coherence and distinctiveness of the final cluster assignments.

% To optimize the merging process and reduce computational overhead, we implement a priority queue that stores cluster pairs sorted by their similarity scores. By always merging the most similar pair first, we ensure that the most semantically coherent clusters are consolidated early in the process. This queue-based approach significantly decreases the time complexity associated with repeatedly searching for the most similar cluster pairs. The algorithm iteratively merges the top pair from the priority queue, updates the TF-IDF vectors and similarity scores for the affected clusters, and re-inserts the updated pairs back into the queue. This process continues until a predefined similarity threshold is met or until the desired number of clusters is achieved.